Read in Gauge Theory Gravity: the metric $g$ below is the derived metric of the gravitational field $(\hb, \Omega)$ on flat spacetime, the manifold is the invertible region, and the tractor connection is the flat-space $\Cl(2,4)$ gauge field of Theorem \claimref{synthesis.gtg-tractor-gauge-field-reproduces-tractor-connection}; the index computation below is its reading in the splitting of $g$ (Fact \claimref{synthesis.imported-gtg-dictionary}).

Over a conformal four-manifold the standard tractor bundle has rank six, carries a metric of signature $(2,4)$, and carries a canonical connection depending only on the conformal class, preserving that metric; its structure group is $SO(2,4)$. Where the conformal manifold is spin, which is assumed throughout (as the correspondence paper assumes it), its spinor bundle is the bundle of conformal spinors (the local twistor bundle), of rank four over the complex unit $I_6$, with fibre $\CS$, on whose fibres the ascended algebra $\Cl(2,4)$ acts with the conformal spinors (twistors) of Fact \claimref{standard-model.imported-pair-systems-two-chirality-halves} as fibre elements; the conformal spinor connection (local twistor transport) is the spinor connection induced by the tractor connection, and its structure group is $\mathrm{Spin}(2,4)$, with Lie algebra $\mathfrak{so}(2,4)$. \emph{Provenance:} the connection formula, \cite[Fact~4.1]{ButlerT6} (\claimref{interior.imported-tractor-connection}); the rank, the signature, the structure group and the spinor bundle, \cite[Fact~4.9, Prop.~4.10]{ButlerT9} (\claimref{correspondence.local-twistor-transport-ascended-spinor-connection}, \claimref{correspondence.imported-tractor-connection-local-twistor-transport}).
